\documentclass{article}
\usepackage[UTF8]{ctex}
\usepackage{graphicx}
\usepackage{fontspec}
\usepackage{xcolor}
\usepackage{amsmath}
\usepackage{amssymb}
\usepackage{bm}
\usepackage{amsfonts} 
\usepackage{tabularx}
\usepackage{booktabs}
\begin{document}
\title{Lecture Notes 8}
\author{Minfei Chen}
\date{\today}
\maketitle
\section{Value Function Approximation: Algorithm}
目标函数：
$$
J(w) = \mathbb{E} \left[ (v_{\pi}(S) - \hat{v}(S, w))^2 \right]
$$

目标是找到$w$，使目标函数，也就是实际的state value与估计的state value相差最小。
\subsection{$S$的概率分布}
$S$为随机变量，它具有一个概率分布。

可能有这些分布形式：
\subsubsection{平均分布}
此时所有状态是同等重要的。

目标函数：
$$
J(w) = \mathbb{E} \left[ (v_{\pi}(S) - \hat{v}(S, w))^2 \right] = \frac{1}{|S|} \sum_{s \in S} (v_{\pi}(s) - \hat{v}(s, w))^2.
$$
\subsubsection{平稳分布 (Stationary Distribution)}
含义: 它表示，如果智能体遵循策略 $\pi$ 与环境进行极长时间的交互，那么在任意一个时间点，我们发现智能体处于状态 s 的概率是多少。

目标函数：
$$
J(w) = \mathbb{E}\left[ (v_{\pi}(S) - \hat{v}(S, w))^2 \right] = \sum_{s \in S} d_{\pi}(s) (v_{\pi}(s) - \hat{v}(s, w))^2.
$$
这个分布给不同的状态赋予了不同的权重。离目标状态更近的状态是更重要的，需要它的估计误差比较小
\section{优化方法}
为了最小化目标函数，需要使用\textbf{梯度下降}的方法：
$$
w_{k+1} = w_k - \alpha_k \nabla_w J(w_k)
$$
\subsubsection*{第一步：计算真实梯度 (The true gradient)}
\begin{align*}
    \nabla_w J(w) &= \nabla_w \mathbb{E}\left[(v_{\pi}(S) - \hat{v}(S, w))^2\right] \\
    % 将梯度算子移入期望内
    &= \mathbb{E}\left[\nabla_w(v_{\pi}(S) - \hat{v}(S, w))^2\right] \\
    % 应用链式法则
    &= \mathbb{E}\left[2(v_{\pi}(S) - \hat{v}(S, w))(-\nabla_w \hat{v}(S, w))\right] \\
    % 整理
    &= -2\mathbb{E}\left[(v_{\pi}(S) - \hat{v}(S, w))\nabla_w \hat{v}(S, w)\right]
\end{align*}

\subsubsection*{第二步：使用随机梯度替代真实梯度}
计算上述包含期望 $\mathbb{E}[\dots]$ 的真实梯度通常是不可行的。因此，我们可以用随机梯度 (stochastic gradient) 来替代它：
\[
w_{t+1} = w_t + \alpha_t(v_{\pi}(s_t) - \hat{v}(s_t, w_t))\nabla_w \hat{v}(s_t, w_t),
\]
其中 $s_t$ 是从状态空间 $S$ 中采样得到的一个样本。这里的 $2\alpha_k$ 被合并简化为了新的 $\alpha_k$。

在求解这个问题时，实际上$v_{\pi}(s_t)$也是未知的。为了得到$v_{\pi}(s_t)$，有两种方法：

方法一：MC learning with function Approximation
$$
w_{t+1} = w_t + \alpha_t (\hat{g}_t - \hat{v}(s_t, w_t))\nabla_w \hat{v}(s_t, w_t).
$$

使用实际中一个episode得到的discounted return $\hat{g}_t$来代替$v_{\pi}(s_t)$。

方法二：TD learning with function Approximation
$$
w_{t+1} = w_t + \alpha_t [r_{t+1} + \gamma \hat{v}(s_{t+1}, w_t) - \hat{v}(s_t, w_t)] \nabla_w \hat{v}(s_t, w_t).
$$
使用TD target $r_{t+1} + \gamma \hat{v}(s_{t+1}, w_t)$来代替$v_{\pi}(s_t)$。
\section{Deep Q-Learning}
深度 Q 学习 (Deep Q-learning) 的目标是最小化以下的目标函数/损失函数：
\[
J(w) = \mathbb{E} \left[ \left( R + \gamma \max_{a' \in A(S')} \hat{q}(S', a', w) - \hat{q}(S, A, w) \right)^2 \right],
\]
其中 $(S, A, R, S')$ 是随机变量。
\subsection*{DQN 的稳定性技巧：固定 Q 目标}

\subsubsection*{问题：移动的目标}

回顾 DQN 的损失函数：
\[
J(w) = \mathbb{E} \left[ \left( R + \gamma \max_{a' \in A(S')} \hat{q}(S', a', w) - \hat{q}(S, A, w) \right)^2 \right],
\]
我们注意到，参数 $w$ 不仅出现在预测值 $\hat{q}(S, A, w)$ 中，也出现在目标值
\[
y \doteq R + \gamma \max_{a' \in A(S')} \hat{q}(S', a', w)
\]
里面。

\begin{itemize}
    \item 为了简化问题，在计算梯度时，我们可以假设目标 $y$ 中的参数 $w$ 是固定的（至少在一段时间内）。
\end{itemize}

\subsubsection*{解决方案一：引入主网络和目标网络}

为了实现上述想法，我们可以引入两个神经网络：
\begin{itemize}
    \item 一个是\textbf{主网络 (main network)}，其参数为 $w$，表示 $\hat{q}(s, a, w)$。
    \item 另一个是\textbf{目标网络 (target network)}，其参数为 $w_T$，表示 $\hat{q}(s, a, w_T)$。
\end{itemize}

在这种情况下，目标函数就退化为：
\[
J(w) = \mathbb{E} \left[ \left( R + \gamma \max_{a' \in A(S')} \hat{q}(S', a', w_T) - \hat{q}(S, A, w) \right)^2 \right],
\]
其中 $w_T$ 是目标网络的参数。
\subsubsection*{解决方案二：经验回放experience replay}
\subsubsection*{理论上的假设}

回顾 DQN 的损失函数：
\[
J(w) = \mathbb{E} \left[ \left( R + \gamma \max_{a' \in A(S')} \hat{q}(S', a', w) - \hat{q}(S, A, w) \right)^2 \right]
\]
在理论分析中，我们通常做出如下假设：
\begin{itemize}
    \item $(S, A) \sim d$：状态-动作对 $(S, A)$ 被视为一个独立的随机变量，遵循某个分布 $d$。
    \item $R \sim p(R|S, A), S' \sim p(S'|S, A)$：奖励 $R$ 和下一个状态 $S'$ 由系统模型决定。
    \item 状态-动作对 $(S, A)$ 的分布被假设为\textbf{均匀分布 (uniform)}。
\end{itemize}

\subsubsection*{现实中的问题与解决方案}

\begin{itemize}
    \item 然而，在实际采集中，样本并不是均匀分布的，因为它们是由某个特定策略按时间顺序生成的。
    \item 为了打破连续样本之间的相关性，我们可以使用经验回放 (experience replay) 技术，即从一个“回放缓冲区 (replay buffer)”中均匀地抽取样本进行训练。
\end{itemize}

这就是为什么我们需要一个回放缓冲区。经验回放也使得样本的利用效率更高（一个样本可以被重复使用）。

\end{document}